% =============================================================
%  Objetivos (Geral e Específicos) (1,00 ponto)
% =============================================================
\section*{Objetivos (Geral e Específicos) (1,00 ponto)}

\textbf{Objetivo geral}

Realizar uma análise comparativa de algoritmos (detecção/rastreamento de objetos, reconhecimento de ações e eventos, detecção de anomalias) que possam complementar um modelo de linguagem no apoio ao trabalho do perito criminal na análise de vídeos de câmeras de monitoramento urbano.

\textbf{Objetivos específicos}

\begin{enumerate}[leftmargin=1.5em]
    \item Levantar na literatura como é composto o trabalho do perito de vídeo (etapas, ferramentas e critérios técnicos hoje empregados), caracterizando o fluxo real de análise a ser apoiado.
    \item Revisar e sistematizar os algoritmos aplicáveis a cada etapa desse fluxo (localização de trechos relevantes, detecção e rastreamento de objetos/pessoas/veículos, reconhecimento de eventos, sumarização e descrição do conteúdo por modelo de linguagem).
    \item Definir critérios de comparação entre os algoritmos/arranjos testados (acurácia, tempo de processamento, exigências computacionais, qualidade e utilidade da análise gerada para fins periciais).
    \item Implementar e testar uma composição de algoritmos com modelo de linguagem sobre um conjunto de vídeos de câmeras de monitoramento urbano, comparando o desempenho das diferentes combinações.
    \item Avaliar a viabilidade prática de uso dessa composição como apoio à rotina de unidades de perícia criminal, discutindo limitações técnicas, éticas e de validade probatória.
\end{enumerate}
